\documentclass[11pt]{article}
\input{style-pc}
\usepackage{array}

\begin{document}

\entetesujet{Sujet bac 2014 -- Physique-Chimie -- Série D}

% ============================ CHIMIE ============================
\matiere{CHIMIE}{8 points}

\partie{A : vérification des connaissances}

\noindent$\blacktriangleright$ \textbf{Appariement.} Relie un élément-question de la colonne A à un élément-réponse de la colonne B. Exemple : $a\cdot8 = b\cdot12$.

\begin{center}\renewcommand{\arraystretch}{1.6}
\begin{tabular}{|p{0.55\linewidth}|p{0.3\linewidth}|}
\hline
\textbf{Colonne A} & \textbf{Colonne B} \\
\hline
$a\cdot1$ Le pH d'une solution de monoacide fort & $b\cdot1$\ \ $t_{\frac{1}{2}} = \dfrac{1}{kC_0}$ \\
\hline
$a\cdot2$ L'élévation ébulliométrique d'une solution & $b\cdot2$\ \ une droite de pente $-k$ \\
\hline
$a\cdot3$ Le temps de demi-réaction d'une solution d'ordre 2 & $b\cdot3$\ \ $-\log C_a$ \\
\hline
$a\cdot4$ Pour une réaction d'ordre 1, la fonction $\ln C = f(t)$ est : & $b\cdot4$\ \ $\Delta\Theta = k'\dfrac{m}{m'M}$ \\
\hline
\end{tabular}
\end{center}

\noindent$\blacktriangleright$ \textbf{Question à réponse courte.} Donne les définitions de :
\begin{enumerate}[label=\alph*.]
  \item Série de raies.
  \item Énergie d'ionisation pour un atome d'hydrogène.
\end{enumerate}

\medskip
\noindent$\blacktriangleright$ \textbf{Réarrangement.} Ordonne le texte suivant qui est écrit en désordre.

\textit{dans l'échantillon soit désintégrée / initialement présents / d'un nucléide / est la durée nécessaire / pour que / la période radioactive / la moitié des noyaux radioactifs}

\partie{B : application des connaissances}
\begin{enumerate}
  \item On prépare une solution $S$ en dissolvant $7{,}9$ g de cristaux anhydres de permanganate de potassium ($\ch{KMnO_4}$) dans 200 mL d'eau. Calcule la concentration molaire volumique de la solution $S$.
  \item On dose en milieu acide 20 mL de la solution $S$ par une solution de sulfate de fer $(\ch{Fe^{2+}} + \ch{SO_4^{2-}})$ à 1 mol/L.
  \begin{enumerate}[label=\alph*.]
    \item Écris l'équation bilan de la réaction d'oxydoréduction.
    \item Calcule le volume de la solution de sulfate de fer utilisé.
    \item Calcule les concentrations molaires volumiques des ions manganèse ($\ch{Mn^{2+}}$) et des ions ferriques ($\ch{Fe^{3+}}$) formés.
  \end{enumerate}
\end{enumerate}
\noindent On donne en g/mol : K = 39 ; Mn = 55 ; O = 16. Couples redox : $\ch{Fe^{3+}/Fe^{2+}}$ ; $\ch{MnO_4^-/Mn^{2+}}$.

% ============================ PHYSIQUE ============================
\matiere{PHYSIQUE}{12 points}

\partie{A : vérification des connaissances}

\noindent\textbf{1. Réponds par vrai ou faux.}
\begin{enumerate}[label=\alph*)]
  \item Un mouvement s'effectuant à une vitesse $v$ constante peut être circulaire ou rectiligne.
  \item Si un solide n'est ni isolé, ni pseudo-isolé, le vecteur accélération de son centre d'inertie n'est pas nul.
  \item Deux vibrations de périodes différentes ne peuvent pas interférer.
  \item L'énergie mécanique d'un système conservatif varie au cours du mouvement.
\end{enumerate}

\medskip
\noindent\textbf{2. Schéma à faire.} Une particule $\alpha$ ($\ch{He^{2+}}$) arrive, avec une vitesse $\vec{V_0}$, en un point $O$ d'un espace champ électrique créé par deux plaques horizontales $A$ et $B$ telles que $V_B - V_A > 0$.

\begin{center}
\begin{tikzpicture}[>=Stealth,scale=1]
  \draw[thick] (0,1)--(5,1) node[right]{$A$};
  \draw[thick] (0,-1)--(5,-1) node[right]{$B$};
  \fill (0.3,0) circle (1.5pt) node[below left]{$O$};
  \draw[->,thick] (0.3,0)--(1.5,0) node[above]{$\vec{V_0}$};
\end{tikzpicture}
\end{center}

\noindent Reproduis puis complète le schéma en représentant le vecteur champ électrique $\vec{E}$ et la force électrique $\vec{F}$ qui s'exerce sur la particule.

\medskip
\noindent\textbf{3. Réarrangement.} Réécris la phrase suivante de manière à définir une grandeur.

\textit{une période / parcourue / la longueur d'onde / la distance / est / par l'onde en}

\partie{B : application des connaissances}
Deux pointes $S_1$ et $S_2$ distantes de 8 cm produisent à la surface horizontale d'une nappe d'eau des vibrations sinusoïdales d'amplitude $a = 2$ mm. Des ondes mécaniques se propagent à la célérité $v = 1{,}5$ m/s.

\begin{enumerate}
  \item La distance entre deux points consécutifs d'amplitude maximale vaut $D = 1$ cm. Détermine la longueur d'onde et la fréquence des vibrations.
  \item Les équations horaires de $S_1$ et $S_2$ sont telles que : $Y_{S_1}(t) = Y_{S_2}(t) = a\sin(2\pi N)t$.
  \begin{enumerate}[label=\alph*.]
    \item Établis l'équation horaire du mouvement d'un point $M$ situé à une distance $d_1$ de $S_1$ et $d_2$ de $S_2$.
    \item Détermine l'élongation de $M$ pour $d_1 = 4$ cm et $d_2 = 7$ cm.
  \end{enumerate}
\end{enumerate}

\partie{C : résolution d'un problème}
En vue de collecter les informations sur un endroit précis du globe terrestre, un satellite doit être placé à une altitude $h$ afin qu'il paraisse immobile pour un observateur terrestre. On dit dans ce cas que ce satellite est géostationnaire.

\begin{enumerate}
  \item Ce satellite, assimilé à un point matériel de masse $m$, doit décrire un mouvement circulaire uniforme à cette altitude $h$. Établis, en fonction de $g_0$, $R$ et $h$ :
  \begin{enumerate}[label=\alph*.]
    \item La vitesse linéaire du satellite.
    \item La période de révolution.
  \end{enumerate}
  \item \begin{enumerate}[label=\alph*.]
    \item Quelle est la valeur de la période de révolution (en secondes) pour que ce satellite soit géostationnaire ?
    \item À quelle altitude $h$ doit-on alors placer ce satellite ?
  \end{enumerate}
\end{enumerate}
\noindent On donne le champ de gravitation terrestre à une altitude $h$ : $g_h = g_0\dfrac{R^2}{(R + h)^2}$ ; $R = 6\,400$ km est le rayon terrestre ; $g_0 = 9{,}8$ m/s$^2$ est le champ de gravitation terrestre au sol.

\end{document}
